\documentclass[10pt,a4paper]{amsart}
\usepackage{amsmath,amssymb,amsthm,graphicx}
\usepackage{verbatim}
\newenvironment{reference}{\comment}{\endcomment}
%\newenvironment{reference}{}{}
\newenvironment{slogan}{\comment}{\endcomment}
\newenvironment{history}{\comment}{\endcomment}

% For commutative diagrams we use Xy-pic
\usepackage[all]{xy}

% We use 2cell for 2-commutative diagrams.
\xyoption{2cell}
% Theorem environments.
%
\theoremstyle{plain}
\newtheorem{theorem}[subsection]{Theorem}
\newtheorem{proposition}[subsection]{Proposition}
\newtheorem{lemma}[subsection]{Lemma}

\theoremstyle{definition}
\newtheorem{definition}[subsection]{Definition}
\newtheorem{example}[subsection]{Example}
\newtheorem{exercise}[subsection]{Exercise}
\newtheorem{situation}[subsection]{Situation}

\theoremstyle{remark}
\newtheorem{remark}[subsection]{Remark}
\newtheorem{remarks}[subsection]{Remarks}

\numberwithin{equation}{subsection}

% Macros
%
\def\lim{\mathop{\mathrm{lim}}\nolimits}
\def\colim{\mathop{\mathrm{colim}}\nolimits}
\def\Spec{\mathop{\mathrm{Spec}}}
\def\Hom{\mathop{\mathrm{Hom}}\nolimits}
\def\Ext{\mathop{\mathrm{Ext}}\nolimits}
\def\SheafHom{\mathop{\mathcal{H}\!\mathit{om}}\nolimits}
\def\SheafExt{\mathop{\mathcal{E}\!\mathit{xt}}\nolimits}
\def\Sch{\mathit{Sch}}
\def\Mor{\mathop{\mathrm{Mor}}\nolimits}
\def\Ob{\mathop{\mathrm{Ob}}\nolimits}
\def\Sh{\mathop{\mathit{Sh}}\nolimits}
\def\NL{\mathop{N\!L}\nolimits}
\def\CH{\mathop{\mathrm{CH}}\nolimits}
\def\proetale{{pro\text{-}\acute{e}tale}}
\def\etale{{\acute{e}tale}}
\def\QCoh{\mathit{QCoh}}
\def\Ker{\mathop{\mathrm{Ker}}}
\def\Im{\mathop{\mathrm{Im}}}
\def\Coker{\mathop{\mathrm{Coker}}}
\def\Coim{\mathop{\mathrm{Coim}}}

% Boxtimes
%
\DeclareMathSymbol{\boxtimes}{\mathbin}{AMSa}{"02}

%
% Macros for moduli stacks/spaces
%
\def\QCohstack{\mathcal{QC}\!\mathit{oh}}
\def\Cohstack{\mathcal{C}\!\mathit{oh}}
\def\Spacesstack{\mathcal{S}\!\mathit{paces}}
\def\Quotfunctor{\mathrm{Quot}}
\def\Hilbfunctor{\mathrm{Hilb}}
\def\Curvesstack{\mathcal{C}\!\mathit{urves}}
\def\Polarizedstack{\mathcal{P}\!\mathit{olarized}}
\def\Complexesstack{\mathcal{C}\!\mathit{omplexes}}
% \Pic is the operator that assigns to X its picard group, usage \Pic(X)
% \Picardstack_{X/B} denotes the Picard stack of X over B
% \Picardfunctor_{X/B} denotes the Picard functor of X over B
\def\Pic{\mathop{\mathrm{Pic}}\nolimits}
\def\Picardstack{\mathcal{P}\!\mathit{ic}}
\def\Picardfunctor{\mathrm{Pic}}
\def\Deformationcategory{\mathcal{D}\!\mathit{ef}}
\usepackage{cleveref}
\crefname{originallemma}{original lemma}{original lemmas}
\Crefname{originallemma}{Original lemma}{Original lemmas}
\makeatletter
\newlabel{lemma-geometrically-irreducible-any-base-change}{{037O}{source}{}{sourcepointer-1}{}}
\newlabel{lemma-geometrically-irreducible-any-base-change@cref}{{[originallemma][0][]037O}{[1][0][]source}{}{}{}}
\makeatother
\crefname{originallemma}{original lemma}{original lemmas}
\Crefname{originallemma}{Original lemma}{Original lemmas}
\makeatletter
\newlabel{lemma-integral-no-inclusion}{{00GT}{source}{}{sourcepointer-2}{}}
\newlabel{lemma-integral-no-inclusion@cref}{{[originallemma][0][]00GT}{[1][0][]source}{}{}{}}
\makeatother
\crefname{originallemma}{original lemma}{original lemmas}
\Crefname{originallemma}{Original lemma}{Original lemmas}
\makeatletter
\newlabel{lemma-make-geometrically-irreducible}{{037Q}{source}{}{sourcepointer-3}{}}
\newlabel{lemma-make-geometrically-irreducible@cref}{{[originallemma][0][]037Q}{[1][0][]source}{}{}{}}
\makeatother
\title{The absolute Galois group acts transitively on the primes of the base change to a separable closure}
\author{The Stacks Project Authors}
\date{stacks-project:algebra.tex:local-snapshot-20261012; GNU FDL 1.2 or later}
\begin{document}
\maketitle
\noindent\textit{Editorial note.} All delivered mathematical text is the author's own, copied from the pinned Stacks Project chapter file algebra.tex; no statement or proof step was rewritten, reordered or removed.\par\medskip
\noindent\textit{Editorial note.} Preamble provenance, as required for this source: the frozen snapshot stages the chapter's own preamble file beside it (preamble.tex, 217 lines, 8116 bytes, the same shared Stacks-wide preamble the published chow and etale-cohomology cards use); the first line of the pinned chapter file is \textbackslash{}input\{preamble\}, so no preamble text lives inside the chapter itself. The card takes that file's lines 8-18 (the verbatim package and the three \textbackslash{}comment-based environments reference/slogan/history at lines 9-12, section 13.2, together with the Xy-pic setup) and its lines 154 to the runtime end of the file (the native theorem declarations and mathematical macros). Not one character was changed, no package was added, no macro was defined, and the excerpt serves only as preamble, never as a span coordinate.\par\medskip
\noindent\textit{Editorial note.} The same-chapter pointers are located at their \textbackslash{}label lines inside the pinned chapter file; the cross-chapter pointers are located in the project's own tag list staged in the snapshot (sources/stacks-chow/tags.txt), at the verbatim lines carrying their canonical tags.\par\medskip
\noindent\textit{Editorial note.} The counted proof refers to two same-chapter statements by label; both are recorded as navigation stubs pointed at their \textbackslash{}label lines, and the arguments consuming them stay inside the counted span.\par\medskip
\noindent\textit{Editorial note.} The counted spans contain no \textbackslash{}cite, so the card carries no bibliography.\par\medskip
\noindent\textit{Allowed standard prerequisites (editorial).}\begin{itemize}
\item The intermediate subextension attached to a field extension. Invoked at the start of the counted proof (lines 11670-11672). Reduces the statement to a subextension over which the base change is geometrically irreducible.
\item Primes of a tensor product over an algebraic extension are maximal. Invoked in the counted proof (lines 11673-11675). Turns the primes whose orbits are studied into maximal ideals, which the Galois action can be compared on.
\end{itemize}\medskip
\ifdefined\hypertarget\hypertarget{sourcepointer-1}{}\fi
\noindent\textit{Source reference.} Original lemma 037O: sources/stacks-algebra/pinned-original-source.tex, source lines 11474--11474. Navigation pointer for \textbackslash{}ref\{lemma-geometrically-irreducible-any-base-change\} used inside the counted span (line 11676); the argument that consumes it is retained in the span.\par\medskip
\ifdefined\hypertarget\hypertarget{sourcepointer-2}{}\fi
\noindent\textit{Source reference.} Original lemma 00GT: sources/stacks-algebra/pinned-original-source.tex, source lines 8030--8030. Navigation pointer for \textbackslash{}ref\{lemma-integral-no-inclusion\} used inside the counted span (line 11675); the argument that consumes it is retained in the span.\par\medskip
\ifdefined\hypertarget\hypertarget{sourcepointer-3}{}\fi
\noindent\textit{Source reference.} Original lemma 037Q: sources/stacks-algebra/pinned-original-source.tex, source lines 11645--11645. Navigation pointer for \textbackslash{}ref\{lemma-make-geometrically-irreducible\} used inside the counted span (line 11672); the argument that consumes it is retained in the span.\par\medskip
\begin{lemma}
\label{lemma-Galois-orbit}
Let $K/k$ be an extension of fields.
Let $\overline{k}/k$ be a separable algebraic closure.
Then $\text{Gal}(\overline{k}/k)$ acts transitively on the
primes of $\overline{k} \otimes_k K$.
\end{lemma}

\begin{proof}
Let $K/k'/k$ be the subextension found in
Lemma \ref{lemma-make-geometrically-irreducible}.
Note that as $k \subset \overline{k}$ is integral all the prime ideals
of $\overline{k} \otimes_k K$ and $\overline{k} \otimes_k k'$ are maximal, see
Lemma \ref{lemma-integral-no-inclusion}.
By Lemma \ref{lemma-geometrically-irreducible-any-base-change}
the map
$$
\Spec(\overline{k} \otimes_k K) \to \Spec(\overline{k} \otimes_k k')
$$
is bijective because (1) all primes are minimal primes, (2)
$\overline{k} \otimes_k K = (\overline{k} \otimes_k k') \otimes_{k'} K$,
and (3) $K$ is geometrically irreducible over $k'$.
Hence it suffices to prove the lemma for the action of
$\text{Gal}(\overline{k}/k)$ on the primes of $\overline{k} \otimes_k k'$.

\medskip\noindent
As every prime of $\overline{k} \otimes_k k'$ is maximal, the residue fields
are isomorphic to $\overline{k}$. Hence the prime ideals of
$\overline{k} \otimes_k k'$ correspond one to one to elements of
$\Hom_k(k', \overline{k})$ with $\sigma \in \Hom_k(k', \overline{k})$
corresponding to the kernel $\mathfrak p_\sigma$ of
$1 \otimes \sigma : \overline{k} \otimes_k k' \to \overline{k}$.
In particular $\text{Gal}(\overline{k}/k)$ acts transitively on
this set as desired.
\end{proof}

\end{document}
